% L13b student companion deck: synthetic markets of many stocks.
% A review of the jump hidden Markov model of one asset with its out-of-sample example, the
% AQR Capital Management company profile, the sector factor model of many stocks (its parts,
% simulating them, the stock-specific draw, the sampling algorithm, and the drifting beta option),
% and the summary.
%
% November 19, 2026 (later): the four computation frames (sector factor, sector loading, stock-specific
% draw, sizing) became one results frame. Their derivations moved to
% advanced/sector_factor/CHEME-5660-L13b-Derivation-SectorFactor-Parts-Fall-2026.ipynb.
% November 19, 2026: the Many Assets frames now tell the story of the lecture: the single index
% model and its covariance, why the single index model is not enough (the first published version
% of the method), the sector factor model, and its block covariance matrix. The Single Index Fit
% frame was removed, because the new opening frame shows the same fit.
% November 18, 2026 (later): the sector factor frames were rewritten to match the lecture's new
% section, which computes the parts from the data first and then fits a single-asset jump hidden
% Markov model to each part.
% November 18, 2026: L13b became the multi-asset lecture only. The trading frames (allocation,
% rebalancing, the market generator, and the INFORMS study's results) left with the trading
% material, parked in lectures/archive/week-13-L13b-trading-2026-11-18/, whose copy of the
% lecture holds the text those frames mirrored. The frame order mirrors the lecture notebook.
% The review, out-of-sample example, and AQR frames were added to match it, and deck-only
% sentences in the factor review and drifting beta frames were replaced with the lecture's
% text. The disclaimer is slide 2, by instructor preference.
% November 17, 2026: the Many Assets frames (sector factor model, drifting beta) moved here
% from L13a.
% Notation follows the lecture: g_i = alpha_i + beta_i g_m + delta_i f_s + psi_i u_i^draw, the
% stock-specific residual u_i and its draw u_i^draw, the drift weight a, the shock size
% k_i, and the standard normal shocks z_i.
\documentclass[aspectratio=169,11pt]{beamer}
\usepackage{vnslides}

\newcommand{\repo}{https://github.com/varnerlab/CHEME-5660-CourseRepository-Fall-2026/blob/main/lectures/week-13/L13b}

\title{\texorpdfstring{{\vnheavy L13b}}{L13b}: Synthetic Markets of Many Stocks}
\subtitle{CHEME 5660 Cornell University}
\author{Copyright Jeffrey Varner 2026. All Rights Reserved}

\begin{document}

\vntitlepage

\begin{frame}{Disclaimer and Risks}
  \vspace{0.6em}
  \textbf{\ink{Training only.}} This content is for training and informational
  purposes only. The teaching team does not make, give, or endorse any offer or
  solicitation to buy or sell securities or derivative products, or provide
  investment or trading advice or strategy.

  \vspace{1.1em}
  \textbf{\ink{Risk.}} Trading involves risk and is generally not recommended for
  those with limited resources, experience, or a low risk tolerance. Past
  performance does not guarantee future success.

  \vspace{1.1em}
  \textbf{\ink{Responsibility.}} You are responsible for your investment decisions
  based on your financial circumstances, objectives, risk tolerance, and
  liquidity needs.
\end{frame}

\begin{frame}{Objectives for Today}
  Previously, we built a jump hidden Markov model of the daily growth rate of one asset.
  How do we simulate many stocks that move together, as real stocks in the same market
  and sector do?

  \vspace{0.5em}
  {\large\textbf{\ink{Today's concepts}}}
  \begin{itemize}\setlength{\itemsep}{0.4em}
    \item \hilite{An out-of-sample test}: paths simulated from the \texttt{SPY} model, with
      and without jump episodes, against the growth rates that came after its fitting period.
    \item \hilite{A synthetic market of many stocks}: a stock-specific baseline plus market,
      sector, and stock-specific parts.
    \item \hilite{Sampling the universe}: a residual scale that keeps each stock's
      variance, with a fixed or drifting market sensitivity.
  \end{itemize}
\end{frame}

\begin{frame}{What You Need to Know}
  \vspace{0.3em}
  {\large\textbf{\ink{Core concepts}}}
  \begin{itemize}\setlength{\itemsep}{0.15em}
    \item An out-of-sample test compares a fitted model with data from after its fitting
      period. We test the \texttt{SPY} model, with and without jump episodes, on the growth
      rates of 2025 and 2026.
    \item We model a synthetic stock's growth rate as
      $g_{i}=\alpha_{i}+\beta_{i}\,g_{m}+\delta_{i}\,f_{s}+\psi_{i}\,u_{i}$.
      The shared market and sector parts make stocks move together.
    \item The residual scale $\psi_i$ matches each stock's variance to the data. A drifting $\beta_i$
      lets the stock's market sensitivity change over time.
  \end{itemize}

  \vspace{0.3em}
  {\large\textbf{\ink{For understanding}} \dimmed{(you won't reproduce these)}}
  \begin{itemize}\setlength{\itemsep}{0.2em}
    \item The tests that compare the \texttt{SPY} model with the out-of-sample data.
    \item How we compute the sector loadings and residual scales, and the drifting $\beta$ option.
    \item The sampling algorithm of the sector factor model, step by step.
  \end{itemize}
\end{frame}

\begin{frame}{Lectures and Examples}
  \vspace{0.3em}
  {\normalsize\textbf{\ink{Lecture:}}\par\small
  \href{\repo/CHEME-5660-L13b-Lecture-JumpHMM-ManyStocks-Fall-2026.ipynb}{\dimmed{CHEME-5660-L13b-Lecture-JumpHMM-ManyStocks-Fall-2026.ipynb}}}

  \vspace{0.6em}
  {\normalsize\textbf{\ink{Examples:}}\par\small
  \href{\repo/CHEME-5660-L13b-Example-JumpHMM-SPY-OutOfSample-Fall-2026.ipynb}{\dimmed{CHEME-5660-L13b-Example-JumpHMM-SPY-OutOfSample-Fall-2026.ipynb}}\par
  \href{\repo/CHEME-5660-L13b-Example-JumpHMM-SectorFactor-Fall-2026.ipynb}{\dimmed{CHEME-5660-L13b-Example-JumpHMM-SectorFactor-Fall-2026.ipynb}}}
\end{frame}

% ========== Review: the jump hidden Markov model of one asset
\begin{frame}{Review: Jump Hidden Markov Models}
  The model bins the daily growth rate into states of equal probability and estimates the
  transition matrix by counting moves between states. Each day, the chain takes one step,
  and the model draws the growth rate from the new state:
  \begin{center}
    \includegraphics[width=\textwidth,height=0.47\textheight,keepaspectratio]{../figs/Fig-L13b-JumpHMM-Schematic.pdf}
  \end{center}
  Jump episodes start at random and last several days. These runs of large moves add the
  volatility clustering that the chain alone misses.
\end{frame}

\begin{frame}{Example: The \texttt{SPY} Model Out of Sample}
  \vspace{0.3em}
  {\normalsize\textbf{\ink{Example:}}\par\small
  \href{\repo/CHEME-5660-L13b-Example-JumpHMM-SPY-OutOfSample-Fall-2026.ipynb}{\dimmed{CHEME-5660-L13b-Example-JumpHMM-SPY-OutOfSample-Fall-2026.ipynb}}}

  \vspace{0.8em}
  In this example, we load the jump hidden Markov model of \texttt{SPY} fitted to daily
  growth rates from 2014 to 2024. We compare its simulated paths with the growth rates of
  2025 and 2026, which the model has not seen. We then turn the jump episodes off, score
  both models on the same data, and compare the volatility clustering of simulated and real
  growth rates.

  \vspace{0.8em}
  \textbf{\ink{Today:}} we extend the jump hidden Markov model from one asset to many stocks
  that move together.
\end{frame}

% ========== Company profile: AQR Capital Management
\begin{frame}{Company Profile: AQR Capital Management}
  \href{https://www.aqr.com/Our-Firm/About-Us}{\hilite{AQR Capital Management}}: a quantitative
  investment management firm, founded in New York City in 1998. It moved its headquarters
  to Greenwich, Connecticut, in 2004.
  \begin{itemize}\setlength{\itemsep}{0.25em}
    \item \hilite{Betting against beta}: in 2014, Andrea Frazzini and Lasse H. Pedersen found
      that stocks with high beta tend to have low alpha. Buying low-beta stocks with borrowed
      money and selling high-beta stocks short earned positive risk-adjusted returns.
  \end{itemize}
  \textbf{\ink{Connection:}} the betting-against-beta portfolio ranks stocks by their market
  sensitivity $\beta_i$, which sets the market part of our sector factor model. With a
  synthetic universe, we could test a betting-against-beta portfolio on many market paths.
  \slidenote{Research:}{\href{https://www.aqr.com/Insights/Research/Journal-Article/Betting-Against-Beta}{Betting Against Beta (2014)}.}
\end{frame}

% ========== Many Assets
\begin{frame}{Many Assets from Single-Asset Models}
  We extend the jump hidden Markov model to many stocks. We start from the single index model,
  then add a sector factor that supplies the covariance the single index model misses.

  \vspace{0.5em}
  We build the many-stock model from \hilite{single-asset jump hidden Markov models}: one for the
  market, one for each sector, and one for each stock.

  \vspace{0.5em}
  Previously, the multiple asset GBM model mixed independent Gaussian shocks through a loading
  matrix. Here, the independent inputs are paths of these single-asset models. Each stock's
  loadings set how much of the market, sector, and stock-specific paths enter its growth rate.
\end{frame}

\begin{frame}{Many Assets: The Single Index Model}
  Previously, we used the single index model to build the covariance matrix of many stocks:
  \begin{equation*}
    g_{i}(t)=\underbrace{\alpha_{i}}_{\textstyle\text{baseline}}+\underbrace{\beta_{i}\,g_{m}(t)}_{\textstyle\text{market}}+\underbrace{\varepsilon_{i}(t)}_{\textstyle\text{single index residual}}
  \end{equation*}
  Here $g_m(t)$ is the growth rate of \texttt{SPY}. The model assumes that the residuals have
  zero mean, are uncorrelated with the market, and are uncorrelated across stocks.
\end{frame}

\begin{frame}{The Single Index Covariance}
  On the diagonal, and off the diagonal for two different stocks $i\neq j$, the second and third
  assumptions give:
  \begin{align*}
    \operatorname{Var}(g_{i}) &= \beta_{i}^{2}\operatorname{Var}(g_{m})+2\beta_{i}\underbrace{\operatorname{Cov}(g_{m},\varepsilon_{i})}_{=0}+\operatorname{Var}(\varepsilon_{i})\\
    \operatorname{Cov}(g_{i},g_{j}) &= \beta_{i}\beta_{j}\operatorname{Var}(g_{m})+\beta_{i}\underbrace{\operatorname{Cov}(g_{m},\varepsilon_{j})}_{=0}+\beta_{j}\underbrace{\operatorname{Cov}(\varepsilon_{i},g_{m})}_{=0}+\underbrace{\operatorname{Cov}(\varepsilon_{i},\varepsilon_{j})}_{=0}
  \end{align*}
  Placing these entries in a matrix gives:
  \begin{equation*}
    \boxed{\symbf{\Sigma}_{g}=\underbrace{\operatorname{Var}(g_{m})\,\boldsymbol{\beta}\boldsymbol{\beta}^{\top}}_{\textstyle\text{market}}+\underbrace{\symbf{\Sigma}_{\varepsilon}}_{\textstyle\text{residual}}}
    \qquad
    \symbf{\Sigma}_{\varepsilon}=\begin{bmatrix}\operatorname{Var}(\varepsilon_{1}) & 0 & \cdots\\ 0 & \operatorname{Var}(\varepsilon_{2}) & \cdots\\ \vdots & \vdots & \ddots\end{bmatrix}
  \end{equation*}
  Every entry has a market part. The residual matrix $\symbf{\Sigma}_{\varepsilon}$ is diagonal,
  because the residuals of two different stocks are uncorrelated.
\end{frame}

\begin{frame}{Why Not the Single Index Model?}
  The first published version of this method used the single index model to simulate 424 stocks
  (\href{https://arxiv.org/abs/2603.10202}{\hilite{Alswaidan and Varner, 2026}}). Every simulated
  stock shared one market draw from the model of \texttt{SPY}, and the residual of each stock was
  drawn independently each day.

  \vspace{0.6em}
  \textbf{\ink{Each stock:}} on the fitting data, the simulated stocks passed the
  Kolmogorov–Smirnov test much less often than the simulated \texttt{SPY} did. The paper traced
  this gap to the residual draws, which did not fully capture the tails of each stock's growth rates.
  We simulate the stock-specific part of each growth rate with its own jump hidden Markov model.

  \vspace{0.6em}
  \textbf{\ink{Stocks in the same sector:}} in the data, the single index residuals of two stocks in
  the same sector usually move together, with $\operatorname{Cov}(\varepsilon_i,\varepsilon_j)>0$.
  The independent residual draws set this covariance to zero. Simulated stocks in the same sector
  then move together less than real stocks do.
\end{frame}

\begin{frame}{Review: Factor Models}
  Previously, we studied \hilite{factor models}, which describe the growth rate in terms of
  factors. For stock $i$ on day $t$, a general factor model is given by:
  \begin{equation*}
    g_{i}(t)=\underbrace{\alpha_{i}}_{\textstyle\text{baseline}}+\underbrace{\beta_{i}\,g_{m}(t)}_{\textstyle\text{market}}+\underbrace{\textstyle\sum_{j}\delta_{ij}\,f_{j}(t)}_{\textstyle\text{other factors}}+\underbrace{\varepsilon_{i}(t)}_{\textstyle\text{residual}}
  \end{equation*}
  Here $\alpha_i$ is the stock-specific baseline (intercept), $g_m(t)$ is the market growth rate, and
  $f_j(t)$ are other factors, such as interest rates or inflation. The coefficients $\beta_i$
  and $\delta_{ij}$ measure sensitivity to these factors. The \hilite{residual}
  $\varepsilon_i(t)$ is the growth left after subtracting the baseline and factor contributions.
\end{frame}

\begin{frame}{Many Assets: The Sector Factor Model}
  The single index model keeps only the market factor. The \hilite{sector factor model} adds one
  factor for each GICS sector, which describes how the single index residuals of its stocks move together:
  \begin{equation*}
    g_{i}(t)=\underbrace{\alpha_{i}}_{\textstyle\text{baseline}}+\underbrace{\beta_{i}\,g_{m}(t)}_{\textstyle\text{market}}+\overbrace{\underbrace{\delta_{i}\,f_{s}(t)}_{\textstyle\text{sector}}+\underbrace{u_{i}(t)}_{\textstyle\text{stock-specific}}}^{\textstyle\text{single index residual }\varepsilon_{i}(t)}
  \end{equation*}
  Here $g_m(t)$ is the growth rate of \texttt{SPY}, $f_s(t)$ is the growth rate of sector $s$
  with the market part removed, and $u_i(t)$ is the growth that neither the market nor the
  sector explains.
\end{frame}

\begin{frame}{The Covariance of the Sector Factor Model}
  The sector factor model splits the single index residual, $\varepsilon_i(t)=\delta_i f_s(t)+u_i(t)$.
  We draw the market, each sector factor, and each stock-specific residual independently. The market
  part of the covariance matrix stays the same, and only the residual matrix $\symbf{\Sigma}_\varepsilon$
  changes. For stock $i$ in sector $s$ and a different stock $j$ in sector $r$, the independent
  draws give:
  \begin{align*}
    \operatorname{Cov}(\varepsilon_{i},\varepsilon_{j}) &= \operatorname{Cov}\left(\delta_{i}\,f_{s}+u_{i},\;\delta_{j}\,f_{r}+u_{j}\right)\\
    &= \delta_{i}\delta_{j}\operatorname{Cov}(f_{s},f_{r})+\delta_{i}\underbrace{\operatorname{Cov}(f_{s},u_{j})}_{=0}+\delta_{j}\underbrace{\operatorname{Cov}(u_{i},f_{r})}_{=0}+\underbrace{\operatorname{Cov}(u_{i},u_{j})}_{=0}\\
    &= \begin{cases}\delta_{i}\delta_{j}\operatorname{Var}(f_{s}) & \text{same sector, } r=s\\ 0 & \text{different sectors, } r\neq s\end{cases}
  \end{align*}
\end{frame}

\begin{frame}{Example: The Sector Blocks}
  On the diagonal, we keep $\operatorname{Var}(\varepsilon_i)$, as in the single index model. Each
  sector adds one block.

  \vspace{0.5em}
  \textbf{\ink{Example:}} Suppose stocks 1 and 2 are in sector $s$, and stocks 3 and 4 are in
  sector $r$. The residual matrix of the four stocks is given by:
  \begin{equation*}
    \symbf{\Sigma}_{\varepsilon}=
    \left[\begin{array}{cc|cc}
    {\color[HTML]{1F6FB4} \operatorname{Var}(\varepsilon_{1})} & {\color[HTML]{1F6FB4} \delta_{1}\delta_{2}\operatorname{Var}(f_{s})} & {\color[HTML]{999999} 0} & {\color[HTML]{999999} 0}\\
    {\color[HTML]{1F6FB4} \delta_{1}\delta_{2}\operatorname{Var}(f_{s})} & {\color[HTML]{1F6FB4} \operatorname{Var}(\varepsilon_{2})} & {\color[HTML]{999999} 0} & {\color[HTML]{999999} 0}\\ \hline
    {\color[HTML]{999999} 0} & {\color[HTML]{999999} 0} & {\color[HTML]{D4720F} \operatorname{Var}(\varepsilon_{3})} & {\color[HTML]{D4720F} \delta_{3}\delta_{4}\operatorname{Var}(f_{r})}\\
    {\color[HTML]{999999} 0} & {\color[HTML]{999999} 0} & {\color[HTML]{D4720F} \delta_{3}\delta_{4}\operatorname{Var}(f_{r})} & {\color[HTML]{D4720F} \operatorname{Var}(\varepsilon_{4})}
    \end{array}\right]
    \begin{array}{l}\left.\vphantom{\begin{matrix}\operatorname{Var}(\varepsilon_{1})\\\operatorname{Var}(\varepsilon_{1})\end{matrix}}\right\}{\color[HTML]{1F6FB4} \text{ sector } s}\\ \left.\vphantom{\begin{matrix}\operatorname{Var}(\varepsilon_{1})\\\operatorname{Var}(\varepsilon_{1})\end{matrix}}\right\}{\color[HTML]{D4720F} \text{ sector } r}\end{array}
  \end{equation*}
  The single index model has the same diagonal and zeros everywhere else.
\end{frame}

\begin{frame}{The Stock-Specific Part}
  The stock-specific draw is a path of daily growth rates simulated from the single-asset jump hidden
  Markov model of stock $i$.
  We multiply the draw by the residual scale $\psi_i$ and add $\psi_i\,u_i(t)$ to the growth rate
  of stock $i$. The sector draw and the
  stock-specific draw are independent.

  \vspace{0.5em}
  We choose the residual scale so that the variance of the simulated single index residual equals
  $\operatorname{Var}(\varepsilon_i)$, the diagonal entry of the residual matrix:
  \begin{equation*}
    \boxed{\psi_i^2=\frac{\operatorname{Var}(\varepsilon_i)-\delta_i^2\operatorname{Var}(f_s)}{\operatorname{Var}(u_i)}}
  \end{equation*}
  The numerator is the variance that the sector part leaves to the stock-specific draw.
\end{frame}

\begin{frame}{Comparing the Two Covariance Matrices}
  In general, each sector adds one block. With the residual scales chosen, the covariance matrices
  of the growth rates of the two models are:
  {\small\begin{align*}
    \text{single index:}\;\symbf{\Sigma}_{g} &= \underbrace{\operatorname{Var}(g_{m})\,\boldsymbol{\beta}\boldsymbol{\beta}^{\top}}_{\textstyle\text{market}} + \underbrace{\operatorname{diag}\big(\operatorname{Var}(\varepsilon_{1}),\dotsc,\operatorname{Var}(\varepsilon_{n})\big)}_{\textstyle\text{residual}}\\[4pt]
    \text{sector factor:}\;\symbf{\Sigma}_{g} &= \underbrace{\operatorname{Var}(g_{m})\,\boldsymbol{\beta}\boldsymbol{\beta}^{\top}}_{\textstyle\text{market}} + \underbrace{\textstyle\sum_{s}\operatorname{Var}(f_{s})\,\boldsymbol{\delta}_{s}\boldsymbol{\delta}_{s}^{\top}}_{\textstyle\text{sector}} + \underbrace{\operatorname{diag}\big(\psi_{1}^{2}\operatorname{Var}(u_{1}),\dotsc,\psi_{n}^{2}\operatorname{Var}(u_{n})\big)}_{\textstyle\text{stock-specific}}
  \end{align*}}%
  \noindent The vector $\boldsymbol{\delta}_s$ holds $\delta_i$ for the stocks in sector $s$ and zeros for the
  others. Both matrices have the same market term. They also have the same diagonal, because
  $\delta_i^2\operatorname{Var}(f_s)+\psi_i^2\operatorname{Var}(u_i)=\operatorname{Var}(\varepsilon_i)$.
  Inside each block, $\delta_i\delta_j\operatorname{Var}(f_s)$ is the residual covariance that the
  single index model sets to zero.
\end{frame}

\begin{frame}{Computing the Parts: The Sector Factor}
  We compute the parts from the data. The sector factor
  is the average single index residual of the $n_s$ stocks in sector $s$:
  \begin{align*}
    f_{s}(t) &= \frac{1}{n_{s}}\sum_{j\in s}\varepsilon_{j}(t)
    = \frac{1}{n_{s}}\sum_{j\in s}\left(g_{j}(t)-\alpha_{j}-\beta_{j}\,g_{m}(t)\right)\\
    &= \frac{1}{n_{s}}\sum_{j\in s}g_{j}(t)-\frac{1}{n_{s}}\sum_{j\in s}\alpha_{j}-\left(\frac{1}{n_{s}}\sum_{j\in s}\beta_{j}\right)g_{m}(t)\\
    &= \underbrace{\bar{g}_{s}(t)}_{\textstyle\text{sector average}}-\underbrace{\bar{\alpha}_{s}}_{\textstyle\text{baseline}}-\underbrace{\bar{\beta}_{s}\,g_{m}(t)}_{\textstyle\text{market}}
  \end{align*}
  Each bar is an average over the stocks in the sector. The sector factor has the market part
  removed. A market move therefore enters the growth rate only once, through $\beta_i g_m(t)$.
\end{frame}

\begin{frame}{Computing the Parts: Loading and Residual}
  \begin{itemize}\setlength{\itemsep}{0.4em}
    \item \hilite{Sector loading}: $\delta_i=c_s\hat{\delta}_i$, where $\hat{\delta}_i$ is the least-squares slope
      of $\varepsilon_i$ on the sector factor without stock $i$. The scale $c_s$ matches the sector block
      to the data on average.
    \item \hilite{Stock-specific residual}: $u_i(t)=\varepsilon_i(t)-\delta_i f_s(t)$ averages zero in the data.
      We then compute the residual scale $\psi_i$.
  \end{itemize}
  The sector factors and the stock-specific residuals are daily growth rates, like the growth rates
  of \texttt{SPY}. We fit a single-asset jump hidden Markov model to each of these growth rate series,
  as we did for \texttt{SPY}.
  \slidenote{Derivations:}{\href{\repo/advanced/sector_factor/CHEME-5660-L13b-Derivation-SectorFactor-Parts-Fall-2026.ipynb}{Computing the parts of the sector factor model}.}
\end{frame}

\begin{frame}{Algorithm: Sample the Synthetic Universe}
  Given the market model (the \texttt{SPY} model tested out of sample), the fitted sector and
  stock-specific models, the coefficients, residual scales, and starting prices:
  \begin{enumerate}\setlength{\itemsep}{0.4em}
    \item \hilite{Shared draws}: draw one market path $g_m(t)$ for all stocks from the
      market model, and one sector path $f_s(t)$ for the stocks in each sector from its
      sector model. Use independent random draws to generate the different paths.
    \item \hilite{Stock-specific draws}: for each stock, draw a path $u_i(t)$ from its
      stock-specific model.
    \item \hilite{Compose}: on each day, calculate
      \begin{equation*}
        g_i(t)=\alpha_i+\beta_i g_m(t)+\delta_i f_s(t)+\psi_i u_i(t),
      \end{equation*}
      then compound $S_i(t)=S_i(t-1)\exp\!\left(g_i(t)\Delta t\right)$.
  \end{enumerate}
  Here the time step $\Delta t=1/252$ year is one trading day.
\end{frame}

\begin{frame}{An Option: A Drifting $\beta$}
  Holding $\beta_i$ fixed assumes that the stock's sensitivity to the market growth rate stays
  the same over a synthetic path. Instead, we let the sensitivity change over
  time around its single index estimate $\hat{\beta}_i$:
  \begin{equation*}
    \beta_i(t)=\hat{\beta}_i+(1-a)\bigl(\beta_i(t-1)-\hat{\beta}_i\bigr)+a k_i z_i(t).
  \end{equation*}
  Here $0<a<1$ is the exponential moving average weight. The independent standard
  normal shocks $z_i(t)$ are scaled by $k_i\geq0$.

  \vspace{0.5em}
  In the example, we choose $k_i$ from how much daily $\beta$ estimates vary in the real data.
  These estimates also contain sampling noise. The market part is now
  $\beta_i(t)g_m(t)$.
  \slidenote{Example:}{\href{\repo/CHEME-5660-L13b-Example-JumpHMM-SectorFactor-Fall-2026.ipynb}{\dimmed{CHEME-5660-L13b-Example-JumpHMM-SectorFactor-Fall-2026.ipynb}}}
\end{frame}

\begin{frame}{Summary}
  \vspace{0.3em}
  We tested the jump hidden Markov model of \texttt{SPY} on data it had not seen. We then
  extended the model from one asset to a universe of stocks that move together.

  \vspace{0.4em}
  {\large\textbf{\ink{Key takeaways}}}
  \begin{itemize}\setlength{\itemsep}{0.4em}
    \item \hilite{Out of sample}: we tested the \texttt{SPY} model on the growth rates of 2025
      and 2026. Paths with a jump episode failed the tests more often and stayed volatile longer.
    \item \hilite{Single-asset models}: we built the universe from single-asset jump hidden Markov
      models for the market, each sector, and each stock. The shared sector draw adds the
      same-sector residual covariance that the single index model sets to zero.
    \item \hilite{Residual scale}: we sized the stock-specific draw so that, with the sector
      draw, it keeps the variance of each stock's single index residual. As an option, we let
      each stock's market sensitivity drift.
  \end{itemize}
  \vspace{0.3em}
  Synthetic markets of many stocks let us test a trading method on many paths, not one.
\end{frame}

\end{document}
